\section{BIẾN CỐ HỢP, BIẾN CỐ GIAO, BIẾN CỐ ĐỘC LẬP}
\subsection{LÝ THUYẾT CẦN NHỚ}

\subsubsection{PHÉP TOÁN TRÊN CÁC BIẾN CỐ}
Cho hai biến cố $A$ và $B$. Khi đó $A, B$ là các tập con của không gian mẫu $\Omega$.
\begin{enumerate}[\iconMT]
	\item \indam{Biến cố hợp:} Biến cố: \lq\lq $A$ hoặc $B$ xảy ra\rq\rq\ được gọi là biến cố hợp của $A$ và $B$, kí hiệu là $A \cup B$.\\
	Biến cố hợp $A \cup B$ là tập con của không gian mẫu $\Omega$.
	\item \indam{Biến cố giao:} Biến cố: \lq\lq Cả $A$ và $B$ đều xảy ra\rq\rq\ được gọi là biến cố giao của $A$ và $B$, kí hiệu là $A B$.\\
	Biến cố giao $A \cap B$ là tập con  của không gian mẫu $\Omega$.
	\item \indam{Biến cố xung khắc:} Biến cố $A$ và biến cố $B$ được gọi là xung khắc nếu $A$ và $B$ không đồng thời xảy ra.
	Hai biến cố $A$ và $B$ xung khắc khi và chỉ khi $A \cap B=\varnothing$.
	\begin{luuy}
		Hai biến cố $A$ và $B$ là xung khắc khi và chỉ khi nếu biến cố này xảy ra thì biến cố kia không xảy ra.
	\end{luuy}
\end{enumerate}
\subsubsection{BIẾN CỐ ĐỘC LẬP}
\begin{enumerate}[\iconMT]
	\item \indam{Định nghĩa:} Cho hai biến cố $A$ và $B$. Hai biến cố $A$ và $B$ được gọi là độc lập nếu việc xảy ra hay không xảy ra của biến cố này không làm ảnh hưởng đến xác suất xảy ra của biến cố kia.
	\item \indam{Chú ý:} Nếu $A, B$ là hai biến cố độc lập thì mỗi cặp biến cố sau cũng độc lập: $A$ và $\overline{B} ; \overline{A}$ và $B ; \overline{A}$ và $\overline{B}$.
\end{enumerate}
\subsubsection{CÁC QUY TẮC TÍNH XÁC SUẤT}
\begin{enumerate}[\iconMT]
	\item \indam{Công thức cộng xác suất:} Cho hai biến cố $A$ và $B$. Khi đó \boxmini{$\mathrm{P}(A \cup B) = \mathrm{P}(A) + \mathrm{P}(B) - \mathrm{P}(A \cap B)$}
	Nếu hai biến cố $A$ và $B$ là xung khắc $\big(A \cap B = \varnothing\big)$ thì 
	\boxmini{$\mathrm{P}(A \cup B)=\mathrm{P}(A)+\mathrm{P}(B)$}
	\item \indam{Công thức nhân xác suất:} Cho hai biến cố $A$ và $B$.
	Nếu hai biến cố $A$ và $B$ là độc lập thì \boxmini{$\mathrm{P}(A B)=\mathrm{P}(A) . \mathrm{P}(B)$}
	\begin{luuy}
		Nếu $\mathrm{P}(A B)\ne \mathrm{P}(A) . \mathrm{P}(B)$ thì $A$ và $B$ không độc lập.
	\end{luuy}
\end{enumerate}

\subsection{PHÂN LOẠI VÀ PHƯƠNG PHÁP GIẢI TOÁN}

\subsection{BÀI TẬP TỰ LUYỆN}

